% Bagean: sajarah
% Bab: biografi
% Pérangan: alan-turing

\documentclass[../../../include/open-logic-section]{subfiles}

\begin{document}

\olfileid{his}{bio}{tur}

\olsection{Alan Turing}

Alan Turing lair ing Maida Vale, London,
tanggal 23 Juni 1912. Dhèwèké dianggep
minangka bapak ilmu komputer téoritis.
Minaté Turing marang ilmu-ilmu fisik
lan matématika wis wiwit enom.
Nanging nalika isih bocah, minaté
ora éntuk papan kang trep ing
sekolah-sekolahé, kang luwih
ngutamakaké sastra lan karya
klasik. Akibaté, bijiné ora
apik lan dhèwèké kerep
ditegur guru-guruné.

\olphoto{turing-alan}{Alan Turing}

Turing kuliah sarjana ing King's
College, Cambridge, lan sinau
matématika. Ing taun 1936
dhèwèké ngembangaké apa kang
saiki diarani mesin Turing
kanggo nyoba ndhéfinisikaké
kanthi cetha gagasan fungsi
kang bisa diitung lan
mbuktèkaké yèn masalah
putusan ora bisa diputusaké.
Alonzo Church wis luwih
dhisik nggayuh asil iku
kanthi kalkulus lambda
duwèké dhéwé. Makalah
Turing tetep kababar
kanthi ngrujuk asil Church.
Church ngundang Turing
menyang Princeton; dhèwèké
ana ing kono saka taun
1936--1938 lan nggayuh
gelar doktor ing sangisoré
bimbingan Church.

Senajan kasengsem marang
logika, minaté Turing
marang ilmu-ilmu fisik
tetep kuwat. Kaprigelan
praktisé digunakaké
nalika dhèwèké ngabdi
ing dhépartemèn
kriptanalisis Inggris
ing Bletchley Park
nalika Perang Donya~II.
Turing dadi tokoh utama
ing mbobol sandi kang
digunakaké komunikasi
angkatan laut Jerman---
sandi Enigma. Kawruhé
Turing ngenani statistika
lan kriptografi, bebarengan
karo panggunaan piranti
èlèktronik, ndadèkaké
timiné bisa mbobol
sandi mau kanthi gawé
mesin pambuka sandi
kang diarani ``bombe.''
Gagasan-gagasané uga
mbantu panggawéyan
komputer èlèktronik
kang bisa diprogram
kapisan ing donya,
Colossus, kang uga
digunakaké ing
Bletchley Park kanggo
mbobol sandi Lorenz
duwèké Jerman.

Turing iku homoseksual.
Senajan mangkono, ing
taun 1942 dhèwèké
nglamar Joan Clarke,
salah siji kanca
sètimé ing Bletchley
Park, nanging banjur
mbatalaké tunangan
lan ngakoni marang
Joan yèn dhèwèké
homoseksual. Sajeroning
uripé dhèwèké
nduwèni sawetara
pacangan, senajan
tumindak homoseksual
wektu iku kalebu
pelanggaran pidana
ing Britania Raya.
Ing taun 1952,
omahé Turing
kémalingan déning
kancané pacangané
wektu iku. Nalika
nglaporaké prakara
iku marang pulisi,
Turing ngakoni
nduwèni sesambungan
homoseksual amarga
ngira pamaréntah
lagi arep nglegalaké
tumindak homoseksual.
Pangira iku salah,
lan dhèwèké didakwa
kanthi tumindak
ora senonoh kang
abot. Tinimbang
dipenjara, Turing
milih pangobatan
hormon kang nyuda
libidoné. Turing
ditemokaké séda
tanggal 8 Juni 1954
amarga overdosis
sianida---kamungkinan
gedhéné amarga bunuh diri.
Dhèwèké éntuk
pangapura karajan
saka Ratu
Elizabeth~II
ing taun 2013.

\begin{reading}
Kanggo biografi Alan Turing
kang jangkep, delengen
\citet{Hodges2014}.
Urip lan karya Turing
dadi inspirasi lakon
\emph{Breaking the Code},
kang digawé dadi
tayangan télévisi
ing taun 1996,
kanthi Derek Jacobi
meranaké Turing.
\emph{The Imitation Game},
film kang dicalonaké
kanggo Academy Award
lan dibintangi
Benedict Cumberbatch
lan Keira Knightley,
uga adhedhasar
kanthi longgar
uripé Alan Turing
lan wektu dhèwèké
ing Bletchley Park
\citep{Imitation2014}.

\citet{Radiolab2012}
nyedhiyakake
sawetara podcast
ngenani urip
lan karya Turing.
Dhokumèntèr
BBC Horizon
\emph{The Strange Life and Death of
  Dr.~Turing}
bisa ditonton
kanthi daring
\citep{Sykes1992}.
\citep{Theelen2012}
iku vidéo cekak
mesin Turing
kang bisa makarya
lan digawé saka
LEGO---kanggo
ngurmati
pengetan 100
taun lairé
Turing ing
taun 2012.

Makalah asli
Turing ngenani
mesin Turing
lan masalah
putusan yaiku
\citet{Turing1937}.
\end{reading}

\end{document}
